\titrePartie{Schémas}

Une règle générale pour les commandes qui permettent de tracer des schémas, c'est que si la commande a pour préfixe tikz, elle doit être utilisée dans un environnement tikzpicture.

Commande pour tracer des rectangles colorés de tailles variables :
\begin{boiteCodeTex}{}
  \rectangle
  \rectangle[largeur = 4cm, hauteur = 1.5cm]
  \rectangle[couleur = couleurSec-300]
  \rectangle[largeur = 8cm, hauteur = 2cm, couleur = couleurTer-400]
\end{boiteCodeTex}

\titreSection{Schémas de planètes, points et vecteurs}

Commande pour tracer des ``planètes'' :
\begin{boiteCodeTex}{}
  \begin{tikzpicture}
    \tikzSchemaPlanete
  \end{tikzpicture}
  \begin{tikzpicture}
    \tikzSchemaPlanete[satellite = B, rayon = 30 pt, orbite = 40 pt, remplissage = couleurQuat-100]
  \end{tikzpicture}
  \begin{tikzpicture}
    \tikzSchemaPlanete[centre = T, satellite = La lune, orbite = 40 pt, contour = couleurSec]
  \end{tikzpicture}
\end{boiteCodeTex}

Quelques commandes pour tracer des points avec un nom et des vecteurs, pour les vecteurs il faut préciser l'origine et la longueur du vecteur selon x et y :
\begin{boiteCodeTex}{cote a cote = 3cm}
  \begin{tikzpicture}
    \tikzLabel (0, 0) {$A$} (-0.1, 0.3)
    \tikzLabel(2, 1) {$B$} (2, 1.3)
    \tikzVecteur[couleurQuat] (0, 0) (2, 1) {$\vv{AB}$} [below = 6pt]
    \tikzVecteur* (0, -0.3) (2, 0) {}
    \tikzLabel* (1, -0.7) {$d$}
  \end{tikzpicture}
\end{boiteCodeTex}

\titreSection{Tube à essais}

Pour tracer des tubes à essais :
\begin{boiteCodeTex}{}
  \begin{tikzpicture}
    \tikzTubeEssais
  \end{tikzpicture}
  %
  \tubeEssaisSolution{couleurTer-400}
  %
  \begin{tikzpicture}
    \tikzPhaseBasTubeEssais [couleur = red-200, hauteur = 0.9]
    \tikzTubeEssais
  \end{tikzpicture}
  %
  \begin{tikzpicture}
    \tikzPhaseBasTubeEssais[hauteur = 0.4, couleur = yellow-200]
    \tikzPhaseTubeEssais[hauteur = 0.4, phase = 0.7, couleur = blue-300]
    \tikzPhaseTubeEssais[hauteur = 1, phase = 1.25, couleur = red-150]
    \tikzTubeEssais
  \end{tikzpicture}
\end{boiteCodeTex}

\titreSection{plan de classe}

Pour faire un plan de classe :
\begin{boiteCodeTex}{}
  \begin{center}
    \begin{tikzpicture}
      \texteRectangle(0, 0) (9, 2) {\large Tableau}
    \end{tikzpicture}
    \bigskip

    \rang (A, B) (C, D, E)
    \rang (F, G) (H) (I, J)
    \rang [largeur = 2, hauteur = 1.5] () () ()
  \end{center}
\end{boiteCodeTex}

%%%%
\begin{boiteCodeTex}{listing only}
  \palette{couleurPrim}{green}
  \palette{couleurSec}{cyan}
  \palette{couleurTer}{blue}
  \palette{couleurQuat}{orange}
\end{boiteCodeTex}
\palette{couleurPrim}{green}
\palette{couleurSec}{cyan}
\palette{couleurTer}{blue}
\palette{couleurQuat}{orange}